%! AP Statistics — Unit 3, Lesson 3.6 — Activity KEY
\documentclass[10pt]{article}
\usepackage{apstats-article}
\usepackage{apstats-key}

\begin{document}

\pageheader{Unit 3, Lesson 3.6}{Group Activity KEY: Choose the Best Design}
\namepartnerperiod

\vspace{0.1in}

\begin{scenariobox}[Context]{sky}
\small
For each research scenario, your group will (1) choose the most appropriate
experimental design (CRD, block design, or matched pairs), (2) justify your
choice, (3) identify any blocking variable, and (4) draw a diagram of the design.
There may be more than one defensible answer --- justify thoroughly.
\end{scenariobox}

\vspace{0.1in}

\textbf{Scenario 1:} A researcher wants to test whether a new math curriculum improves
student test scores. She has 60 students who are evenly split between students who
are grade-level and students who are below grade-level in math.

\begin{practicebox}
\small
\begin{enumerate}[leftmargin=1.5em, itemsep=8pt]
  \item Recommended design: \makebox[6cm][l]{\ans{Block design}}
  \item Blocking variable (if any): \makebox[6cm][l]{\ans{Math level (grade-level vs.\ below)}}
  \item Justify your choice: \makebox[9cm][l]{\ans{Math level likely affects test scores, so}}\\
        \ansline{blocking removes that variability and sharpens the comparison.}
  \item \makebox[0pt][l]{\raisebox{-2\baselineskip}[0pt][0pt]{\parbox[t]{0.9\linewidth}{\ans{60 students $\to$ Block 1: grade-level (30) and Block 2: below grade-level (30).\\
        Within each block, randomly assign 15 to the new curriculum and 15 to the\\
        traditional curriculum $\to$ compare test scores within each block.}}}}Draw a diagram of your design:
  \vspace{1.4in}
\end{enumerate}
\end{practicebox}

\vspace{0.1in}

\textbf{Scenario 2:} A dermatologist wants to compare two acne treatments (Gel A
vs.\ Gel B). She has 40 patients; each has two clearly affected areas (left cheek
and right cheek) with equal severity.

\begin{practicebox}
\small
\begin{enumerate}[leftmargin=1.5em, itemsep=8pt]
  \item Recommended design: \makebox[6cm][l]{\ans{Matched pairs}}
  \item Explain why this design is ideal here: \makebox[9cm][l]{\ans{Each patient is their own pair: both cheeks}}\\
        \ansline{share genetics, diet, and hygiene, so those are controlled.}
  \item How would you assign treatments to left vs.\ right cheek?\\
        \ansline{For each patient, flip a coin: heads $\to$ Gel A on the left cheek,}\\[1pt]\ansline{Gel B on the right; tails $\to$ the reverse.}
  \item \makebox[0pt][l]{\raisebox{-2\baselineskip}[0pt][0pt]{\parbox[t]{0.9\linewidth}{\ans{40 patients $\to$ for each patient, coin flip assigns Gel A to one cheek and\\
        Gel B to the other $\to$ compare improvement (A $-$ B) within each patient.}}}}Draw a diagram:
  \vspace{1.2in}
\end{enumerate}
\end{practicebox}

\vspace{0.1in}

\textbf{Scenario 3:} A seed company tests three types of corn seed (A, B, C) on
48 plots. The plots are spread across 4 fields with different soil quality.
There is no way to use the same plot for multiple treatments.

\begin{practicebox}
\small
\begin{enumerate}[leftmargin=1.5em, itemsep=8pt]
  \item Recommended design: \makebox[6cm][l]{\ans{Block design}}
  \item Blocking variable: \makebox[6cm][l]{\ans{Field (soil quality)}}
  \item How many plots per treatment per block? \makebox[3cm][l]{\ans{4}}
  \item \makebox[0pt][l]{\raisebox{-2\baselineskip}[0pt][0pt]{\parbox[t]{0.9\linewidth}{\ans{4 blocks (fields) of 12 plots each. Within each field, randomly assign\\
        4 plots to Seed A, 4 to Seed B, 4 to Seed C $\to$ compare yields within each field.}}}}Draw a diagram showing all 48 plots:
  \vspace{1.4in}
  \item One advantage of this design over a CRD: \makebox[9cm][l]{\ans{Removes field-to-field soil variability.}}
\end{enumerate}
\end{practicebox}

\vspace{0.12in}

\begin{practicebox}
\small
\textbf{Decision Chart:} Complete the chart your group can use on the AP exam.

\vspace{4pt}
\renewcommand{\arraystretch}{1.8}
\begin{tabularx}{\linewidth}{@{} X c @{}}
\textbf{Situation} & \textbf{Best design} \\
\hline
Units are homogeneous; no known nuisance variable. & \ans{CRD} \\
A known variable affects the response; units cannot be paired. & \ans{Block design} \\
Units can be paired 1-to-1 on a key variable. & \ans{Matched pairs} \\
Same unit can receive both treatments in sequence. & \ans{Matched pairs} \\
\end{tabularx}
\end{practicebox}

\end{document}
